% \begin{figure}[tb]
% \centering
% \begingroup
% \tikzset{
%   overviewfile/.style={draw, rounded corners=1pt, fill=gray!15, minimum width=8mm, minimum height=5mm, font=\scriptsize},
%   overviewmember/.style={draw, circle, inner sep=2pt, font=\scriptsize},
%   overviewfolder/.style={draw, fill=gray!35, minimum width=12mm, minimum height=5mm, font=\scriptsize},
%   overviewedge/.style={->, thin}
% }
% \newcommand{\variantpanel}[3]{%
% \begin{tikzpicture}[x=1cm,y=1cm,baseline=(current bounding box.center)]
%   \path[use as bounding box] (-0.55,-0.8) rectangle (2.15,1.25);
%   \node[font=\small\bfseries] at (0.8,1.08) {#1};
%   \node[overviewfile] (f1) at (0,0) {$f_1$};
%   \node[overviewfile] (f2) at (1.6,0) {$f_2$};
%   \draw[overviewedge] (f1) -- (f2);
%   \ifnum#2=1
%     \node[overviewfolder] (d) at (0.8,0.65) {Folder};
%     \draw[overviewedge,dashed] (d) -- (f1);
%     \draw[overviewedge,dashed] (d) -- (f2);
%   \fi
%   \ifnum#3=1
%     \node[overviewmember] (m1) at (0,-0.55) {$m_1$};
%     \node[overviewmember] (m2) at (1.6,-0.55) {$m_2$};
%     \draw[overviewedge,dashed] (m1) -- (f1);
%     \draw[overviewedge,dashed] (m2) -- (f2);
%     \draw[overviewedge] (m1) -- (m2);
%   \fi
% \end{tikzpicture}%
% }
% \begin{tabular}{@{}c@{\hspace{6mm}}c@{}}
% \variantpanel{F}{0}{0} & \variantpanel{FM}{0}{1} \\
% \variantpanel{FF}{1}{0} & \variantpanel{FFM}{1}{1}
% \end{tabular}
% \endgroup

\begin{figure}[tb]
\centering
\begingroup
\tikzset{
  overviewfile/.style={
    draw,
    circle,
    fill=gray!15,
    minimum size=7mm,
    inner sep=1pt,
    font=\scriptsize
  },
  overviewmember/.style={
    draw,
    circle,
    inner sep=1.5pt,
    font=\scriptsize
  },
  overviewfolder/.style={
    draw,
    fill=gray!35,
    minimum width=12mm,
    minimum height=6mm,
    font=\scriptsize
  },
  overviewedge/.style={
    -{Stealth[length=1.8mm,width=1.3mm]},
    line width=0.6pt
  }
}
\newcommand{\variantpanel}[3]{%
\begin{tikzpicture}[
  x=1cm,
  y=1cm,
  baseline=(current bounding box.center)
]
  % Reserve folder space only for FF and FFM.
  \ifnum#2=1
    \path[use as bounding box]
      (-0.55,-1.2) rectangle (2.55,1.65);
    \node[font=\small\bfseries] at (1,1.45) {#1};
  \else
    \path[use as bounding box]
      (-0.55,-1.2) rectangle (2.55,0.8);
    \node[font=\small\bfseries] at (1,0.6) {#1};
  \fi

  \node[overviewfile] (f1) at (0,0) {$f_1$};
  \node[overviewfile] (f2) at (2,0) {$f_2$};
  \draw[overviewedge] (f1.east) -- (f2.west);

  \ifnum#2=1
    \node[overviewfolder] (d) at (1,0.95) {Folder};
    \draw[overviewedge,dashed]
      (d.south west) -- (f1.north);
    \draw[overviewedge,dashed]
      (d.south east) -- (f2.north);
  \fi

  \ifnum#3=1
    \node[overviewmember] (m1) at (0,-0.9) {$m_1$};
    \node[overviewmember] (m2) at (2,-0.9) {$m_2$};
    \draw[overviewedge,dashed]
      (m1.north) -- (f1.south);
    \draw[overviewedge,dashed]
      (m2.north) -- (f2.south);
    \draw[overviewedge]
      (m1.east) -- (m2.west);
  \fi
\end{tikzpicture}%
}

\begin{tabular}{@{}c@{\hspace{6mm}}c@{}}
  \variantpanel{F}{0}{0}
  & \variantpanel{FM}{0}{1} \\
  \variantpanel{FF}{1}{0}
  & \variantpanel{FFM}{1}{1}
\end{tabular}
\endgroup

\caption{Illustrative graph configurations. Files $f_i$ are the prediction targets; each $m_i$ aggregates the members of file $f_i$. Solid arrows show example dependencies and dashed arrows show containment. Moving right adds member nodes; moving down adds folders. The panels omit relation types, additional dependency edges and folder ancestry; Figure~\ref{fig:schema} gives the typed schema. Feature construction also depends on the configuration, as described in Section~\ref{sec:node_features}.\label{fig:graph_variants_overview}}
\end{figure}
